We start by computing the f\/irst term in the right hand side of formula
\eqref{quasir}, which is $\dim g^1_u(P_1)$, for some $u\in
TM\setminus\{0\}$. Recall that $g^1(P_1)=\operatorname{Ker} \left(\sigma^1(P_1)\right)
\subset T^*\otimes T^*_v$.  We have to compute the dimension of
the f\/ibers of $g^1(P_1)$, which is a vector subbundle of $ T^*\otimes T^*_v$. An element $A\in
g^1(P_1)$ can be expressed, with respect to the adapted dual basis
$\{dx^i, \delta y^i\}$, as follows
\begin{gather*}
A=A_{ij}dx^i\otimes dx^j + A_{\underline{i}j}\delta y^i \otimes
dx^j. %\label{aij}
\end{gather*} Using formula \eqref{sigma1p1},
the symbol $\sigma^1(P_1)$ can be expressed as follows:
\[
\sigma^1(P_1)A=\left( A_{\underline{i}j}y^idx^j, \frac{1}{2}(
A_{\underline{i}j}- A_{\underline{j}i})dx^i\wedge dx^j,
\frac{1}{2}(A_{ij}- A_{ji})dx^i\wedge dx^j\right).
\]
The condition $\tau_hA=0$ is equivalent with $A_{ij}=A_{ji}$ and due
to this condition $A_{ij}$ contribute with $n(n+1)/2$ to the $\dim g^1_u(P_1)$. The conditions
$\tau_JA=0$ and $i_{\mathbb{C}}A=0$ are equivalent to
$A_{\underline{i}j}=A_{\underline{j}i}$, and respectively
$A_{\underline{i}j}y^i=0$. Hence, due to these two conditions,
$A_{\underline{i}j}$ contribute
with $n(n-1)/2$ to the the $\dim g^1_u(P_1)$. It follows that $\dim
g^{1}_u(P_1)=n(n-1)/2 + n(n+1)/2 = n^2.$

We continue the proof by computing the left hand side of formula
\eqref{quasir}, which is $\dim g^2_u(P_1)$.
Therefore, we will consider the kernel of the f\/irst-order prolongation of the symbol,
$g^2(P_1)=\operatorname{Ker}\left(\sigma^2(P_1)\right)\subset
S^2T^*\otimes T^*_v$. An element $B\in g^2(P_1)$ can be expressed,
with respect to the adapted dual basis $\{dx^i, \delta y^i\}$, as
follows
\begin{gather*}
B=B_{ijk}dx^i\otimes dx^j\otimes dx^k + B_{\underline{i}jk}\delta
y^i \otimes dx^j \otimes dx^k \\
\phantom{B=}{} + B_{i\underline{j}k}dx^i \otimes
\delta y^j \otimes dx^k + B_{\underline{ij}k}\delta y^i \otimes
\delta y^j \otimes dx^k,
\end{gather*}
with the symmetry conditions $B_{ijk}=B_{jik}$,
$B_{\underline{i}jk}=B_{i\underline{j}k}$ and
$B_{\underline{ij}k}=B_{\underline{ji}k}$ satisf\/ied. Using formula~\eqref{sigma2p1}, the symbol $\sigma^2(P_1)$ can be expressed as
follows
\begin{gather*}
\sigma^2(P_1)B  = \bigg(B_{i\underline{j}k}y^j dx^i \otimes dx^k +
B_{\underline{ij}k}y^j \delta y^i \otimes dx^k,  \\
\hphantom{\sigma^2(P_1)B  = \bigg(}{}  \frac{1}{2}(B_{i\underline{j}k} -
B_{i\underline{k}j})dx^i\otimes dx^j \wedge dx^k +
\frac{1}{2}(B_{\underline{ij}k} - B_{\underline{ik}j})\delta
y^i\otimes dx^j \wedge dx^k, \\
\hphantom{\sigma^2(P_1)B  = \bigg(}{}  \frac{1}{2}(B_{ijk} - B_{ikj})dx^i\otimes dx^j \wedge
dx^k + \frac{1}{2}(B_{\underline{i}jk} -
B_{\underline{i}kj})\delta y^i\otimes dx^j \wedge dx^k\bigg).
\end{gather*}
The totally symmetric components $B_{ijk}$ contribute with
$n(n+1)(n+2)/6$ to the $\dim g^2_u(P_1)$. The other two are also
totally symmetric components on the $(n-1)$-dimensional space
given by restrictions $B_{i\underline{j}k}y^j=0$ and respectively
$B_{\underline{ij}k}y^j=0$. Therefore, each of them contributes
with $(n-1)n(n+1)/6$ to the $\dim g^2_u(P_1)$. Consequently, $\dim
g^2_u(P_1)=n(n+1)(n+2)/6+2(n-1)n(n+1)/6=n^2(n+1)/2$.
